\section{Glacial Recovery across the I/O Barrier}
\label{sec:recovery}

Recovery has two obligations: read the deployment value out of the configuration file, and prove that treating it as a constant is sound. The second is the crux, and it is what gives this step content beyond a neck-split that simply trusts its split point.

\subsection{Recovering the value}
Each \lstinline[style=fortran]|READ(NML=g)| statement names a namelist group \texttt{g} whose members are module-level variables, or fields of a configuration derived type. We map each member to the storage it initializes and read its value from the deployment namelist file; a variable a group reads is recovered as the value the file gives, or, across several deployment files, as the small set of values they give (\Cref{sec:variants}). Recovery is transitive: a derived variable is glacial when everything it is computed from is glacial and the computation lies in the init phase. The LU copy $\mathit{nx}=\mathit{nx0}$ is the canonical case --- $\mathit{nx}$ is recovered because $\mathit{nx0}$ is and \lstinline[style=fortran]|domain()| runs before the timestep loop --- and it is also why the derived bound, not the namelist variable, is the one the compiler leaves as a runtime load (\Cref{tab:lu}).

\subsection{The glacial condition}
Injecting a recovered value is sound only if the variable does not change during the compute phase that reads it. We certify this with a may-write analysis: a variable is \emph{glacial}, hence recoverable, when no store to its storage is reachable from any timestep-loop entry. The obligation is only as strong as the set of write channels the analysis covers, and production Fortran has many: direct assignment, \texttt{COMMON}/\texttt{EQUIVALENCE} aliasing, pointer association, \texttt{INTENT(OUT)} and \texttt{INTENT(INOUT)} arguments, stores through \texttt{c\_loc}/\texttt{c\_f\_pointer}, writes through procedure pointers and runtime-resolved generic interfaces, and, in a distributed run, the receive buffers of MPI and one-sided RMA calls that may alias configuration storage. The analysis must account for all of them; where it cannot prove the write set empty, it \emph{refuses}, and the variable stays generic. Refusal is the fail-safe: a missed channel costs a specialization, never correctness.

One dependency is circular. A sound whole-program may-write set needs the targets of indirect dispatch, but resolving that dispatch is exactly the devirtualization of \Cref{sec:spec}, which itself consumes recovered values. We break the cycle by over-approximating indirect targets conservatively during recovery --- every type-compatible procedure is a possible callee --- and devirtualizing only afterward. The conservative step can only add writes, so it can only make the analysis refuse, never wrongly accept.

\subsection{Recovery as automatic staging}
A configured code is a two-stage program: the init phase consumes the configuration and writes the glacial state, and the timestep loop reads it. Recovery is the binding-time analysis that finds the boundary between the stages and executes the first at deployment-compile time. The boundary is not the textual \lstinline[style=fortran]|READ|: a value is often finished later, after configuration derivation, an MPI broadcast to the ranks, resolution-dependent re-defaulting, or a restart override. We therefore define it per variable as the last program point that dominates every timestep-loop entry and past which the variable's compute-phase may-write set is empty. This locates staging automatically, where classical partial evaluation and SFAC take the binding-time split by hand, and Autrey and Wolfe~\cite{glacial} describe the staging without acting on it.

\subsection{Positioning}
LMCAS~\cite{lmcas} splits a C program at a neck between argument parsing and main logic and constant-propagates the captured state; SFAC~\cite{sfac} specializes on hand-supplied values. Neither recovers persisted module-level state across an I/O barrier in multi-translation-unit Fortran, and neither carries the glacial certification above, which is what makes the recovered value safe to treat as a constant.

\td{OPEN: prove the per-variable staging boundary is well-defined (a unique last dominating point with empty compute-phase may-write set exists), or admit per-variable boundaries and show residualization is still sound. Tie the channel list to what the implementation actually checks; mark any channel currently handled by refusal rather than analysis.}
